\documentclass[letterpaper]{article}

\usepackage[margin=1.5cm]{geometry}
\usepackage{amsmath,amsfonts}
\usepackage{url}
\usepackage{bm}

\newcommand{\Real}{\mathbb{R}}
\newcommand{\tr}{^\top}
\newcommand{\opt}{^\star}

\setlength{\parskip}{3mm plus1mm minus1mm}
\setlength{\parindent}{0mm}

\title{Problem Set 12}
\author{Foundations of Machine Learning}
\date{}

\begin{document}
\maketitle

\begin{itemize}
\item \textbf{Submission}: Turn in as a \texttt{PDF} and a \textbf{separate source code file} on \url{http://mycourses.unh.edu}. Scanned hand-written notes are encouraged.
\item \textbf{Questions}: See the class website for a link to office hours
\end{itemize}

% Covers: Maximum likelihood and MAP (MML 8.1, 8.3) and Principal component analysis (MML 10.1--10.3).

\textbf{Relevant reading}: MML 8.1, 8.3 (maximum likelihood, MAP); MML 10.1--10.3 (principal component analysis). Problems marked \textbf{[code]} should be submitted as a program in a language of your choice (Julia or Python recommended). The cited exercise numbers are from MML (Deisenroth et al., 2021).

\begin{enumerate}
\item MML, Chapter 8 exercises: any two on maximum likelihood and MAP estimation (e.g.\ Exercises 8.1--8.3).
\item MML, Chapter 10 exercises: any two on principal component analysis (e.g.\ Exercises 10.1--10.3).
\item \textbf{[code]} Given a sample of data, compute the maximum-likelihood estimates of the mean and variance of a Gaussian and of the parameter of a Bernoulli random variable, and compare each with the MAP estimate under a simple prior.
\item \textbf{[code]} Run PCA on a small data set: center the data, form the covariance matrix, and compute its principal components. Plot the fraction of variance explained by each component and project the data onto the first two components.
\item \textbf{Time spent.} Report approximately how many hours you spent completing this problem set (including the reading).
\end{enumerate}

\end{document}
